\chapter{The run session: what a layout binds to}\label{ch:session}

\texttt{HeadlessRunSession} in \texttt{src/\allowbreak{}ate/\allowbreak{}run\_\allowbreak{}session/\allowbreak{}headless.py} is the object a screen talks
to, and the object a headless run uses. It is a \texttt{QObject}: that is what lets a widget connect to
its signals. It imports no widgets, and it writes no SQL. It holds the Operator's decisions, owns one
engine at a time, and turns the worker's signals into its own.

\section{The states}

\texttt{SessionState} has four values, and the session keeps \texttt{\_state} in step with what it
emits: \texttt{IDLE} (nothing running, fields not yet entered), \texttt{SELECTING} (a schedule is
selected, the run has not started), \texttt{STARTED} (the worker is running), \texttt{FINISHED}. The
engine's \texttt{RunStatus} is mapped onto these by \texttt{\_map\_run\_status}: \texttt{RUNNING}
becomes \texttt{STARTED}.

\texttt{STARTED} is announced on the calling thread before the worker starts, so a screen can paint the
run view before the first signal arrives. The worker's queued signals then move the state.

\section{The nine signals}

\begin{itemize}
  \item \texttt{state\_changed}, \texttt{input\_asked}, \texttt{step\_finished}, \texttt{failure\_prompt},
    \texttt{prompt\_answered}, \texttt{action\_started}, \texttt{step\_started}: forwarded from the
    engine, so a layout connects to the session and never to the engine.
  \item \texttt{input\_answered} (node seq, Attempt no, index, value): the session's own, emitted when an
    answer it accepted reached the worker. Its value type is \texttt{object}, because a Step Input can be
    a number, a string or a list of labels.
  \item \texttt{report\_written} (HTML path, PDF path): fired exactly once per finished run, with an
    empty second path when there is no PDF --- because the render was skipped, or because it failed.
\end{itemize}

That last one is worth the space it takes. The report is written after \texttt{FINISHED} is announced,
and the PDF is rendered by a \texttt{QProcess} the session owns (chapter~\ref{ch:pdf}). The sandbox
question is asked \emph{after} the HTML is on disk, so a declined PDF still leaves the report and the
signal still fires once. A session that never gets a \texttt{report\_written} is a session whose teardown
was wrong.

\section{Rules the session enforces}

The session, not the screen, is the authority:

\begin{itemize}
  \item \texttt{start()} refuses to run with no schedule selected, and refuses a run with an empty
    Project (spec 08/11). A screen that disables the Start button is being helpful; the session is what
    makes it true.
  \item Before a run starts, \texttt{close\_\allowbreak{}interrupted\_\allowbreak{}runs} closes every run row with no
    \texttt{ended\_at} --- one left open by a crash or a power cut --- recording it
    \texttt{INTERRUPTED}/\texttt{FAIL}, and logs how many. That is spec rule 98, and it is why a run
    never stays open across launches.
  \item Stop on Fail is decided in one place: the Operator's explicit choice wins, otherwise the
    schedule's own default (captured when the schedule was selected, \texttt{True} when unset). Whatever
    it resolves to is what the run row records.
  \item \texttt{supply\_input} and \texttt{click\_poll} check the answer before the worker sees it: the
    step and Attempt must be the ones running, the index must be declared by that step, the input must
    actually be waiting, and the value must satisfy the declaration (packet 10a review). A bad value
    raises \texttt{ValueError} and is never recorded. \texttt{index} here is the 1-based \texttt{n} of
    \texttt{ATE.input(n)}; \texttt{OperatorInput} keys its channels from 0, and this is where the two
    meet.
  \item \texttt{pause\_pending()} is how a layout knows the run is paused. The layouts use the session,
    not their own prompt bar, because the bar of the layout that is not showing says nothing.
\end{itemize}

\section{Closing, and why the order matters}

\texttt{close()} ends the worker, writes the HTML report if it has not been written, marks the session
closed so a still-queued \texttt{FINISHED} finds nothing to write, stops the PDF helper, closes the
store, and closes the application log only if the session opened its own (a log handed in by the launch
belongs to the launch). \texttt{\_\_del\_\_} does the helper stop, because a session dropped without
\texttt{close()} leaves its \texttt{QProcess} to be destroyed by Qt mid-render, and the job's
\texttt{finished} callback then fires against a session whose Qt side is gone (packet 25b review: this
was the teardown crash).

Do not close a store while a worker is still writing to it. \texttt{close()} calls
\texttt{stop\_and\_wait()} first for exactly that reason; if the worker does not stop within 5 s it logs
an error rather than pretending it did.

\section{Driving a run with no display}

This is the shape the suite and the harness use, and it is the shape to copy for a new headless test.
The runner below is \texttt{docs/manuals/examples/developer/runner.py}, the file the manual's worked
example in chapter~\ref{ch:device-type} runs; \texttt{tests/test_manuals.py} runs it as a subprocess.

\lstinputlisting[language=ATEPython, caption={\texttt{docs/manuals/examples/developer/runner.py}}]{../examples/developer/runner.py}

There is no widget and no display, and the queued signals still deliver, because delivery rides the
event loop of the receiving thread. Without the \texttt{processEvents} pump nothing arrives: no
\texttt{state\_changed}, no step recorded, no report. A Stop-on-Fail pause blocks the worker until
somebody answers it, and in a headless run the script is the only somebody, which is why the runner
answers every pause from \texttt{--choice}. A blocking Step Input is answered the same way, from
\texttt{--input STEP NAME=VALUE}: \texttt{input\_asked} says which node is waiting, the Step Name is
looked up in the schedule's tree, and \texttt{supply\_input} hands the value over. The script exits
0 only when the session reached \texttt{FINISHED}.

\cando{start a Test Run from a script, answer its Step Inputs, abort it, and read the record back from
the store.}
